\chapter{Test} \label{ch:Test}

\newpage
\Nref{ch:References}
\autocite{TeXtured}
If you encounter any issues, please see \href{https://github.com/jdujava/TeXtured/issues/2}{\texttt{jdujava/TeXtured \#2}} and \href{https://github.com/jdujava/TeXtured/issues/5}{\texttt{jdujava/TeXtured \#5}}.


% 1. Inline usage
This is a sentence about the X-Band pipeline \todo{Verify bandwidth parameters}. 
You can also categorize them: \todo[Urgent]{Fix the Viterbi rate calculation}.

% 2. Full environment usage
\begin{Todo}[Missing References]
    We need to add the specific Blue Book versions for the AOS Space Data Link Protocol here.
\end{Todo}

\begin{Question}
    Should we use a 1/2 or 3/4 convolutional rate for the Astra benchmarking?
\end{Question}

\begin{Suggestion}[Optimization]
    Consider replacing the standard PLL with an EKF to improve phase tracking in high-Doppler LEO scenarios.
\end{Suggestion}

\begin{Note}
    The current LNB has a TCXO with 1 ppm stability, which should be sufficient for initial X-Band tests.
\end{Note}

H00pa \\\textcolor{gray}{at first subconscious}


\autocite[pág. 20]{CCSDS_131_0_B_5}

\newpage
De acuerdo con el artículo científico, los autores caracterizan el efecto Doppler utilizando el término \textbf{Doppler normalizado} y derivan expresiones analíticas aproximadas simplificando la geometría orbital.

\section{Expresión del Efecto Doppler (Doppler Normalizado)}

La expresión analítica exacta del Doppler normalizado, definida en el artículo como la relación entre el cambio de frecuencia y la frecuencia transmitida ($\frac{\Delta f}{f}$), corresponde a la \textbf{ecuación (5)}:

\begin{equation}
\frac{\Delta f}{f} = -\frac{1}{c}\frac{r_{E}r \sin(\psi(t)-\psi(t_{0}))\cos\left(\cos^{-1}\left(\frac{r_{E}}{r}\cos \theta_{max}\right)-\theta_{max}\right)\cdot\omega_{F}(t)}{\sqrt{r_{E}^{2}+r^{2}-2r_{E}r \cos(\psi(t)-\psi(t_{0}))\cos\left(\cos^{-1}\left(\frac{r_{E}}{r}\cos \theta_{max}\right)-\theta_{max}\right)}}
\end{equation}

\subsection*{Nomenclatura de variables:}
\begin{itemize}
    \item $c$: Velocidad de la luz en el vacío.
    \item $r_{E}$: Radio de la Tierra.
    \item $r$: Distancia desde el centro de la Tierra al satélite ($r = r_{E} + h$, donde $h$ es la altitud de la órbita).
    \item $\theta_{max}$: Ángulo de elevación máximo observado por la terminal terrestre durante la ventana de visibilidad.
    \item $t_{0}$: Instante de tiempo en el que ocurre el Doppler cero (coincide con el momento de máxima elevación).
    \item $\psi(t)-\psi(t_{0})$: Distancia angular recorrida por la traza del satélite sobre la superficie terrestre desde el instante de referencia $t_{0}$.
    \item $\omega_{F}(t)$: Velocidad angular del satélite medida en el sistema de referencia de la Tierra Fija (ECF).
\end{itemize}

\section{Expresión de la Curva (Aproximación de Velocidad Angular Constante)}

Para modelar matemáticamente la curva en forma de ``S'' que describe el comportamiento del Doppler a lo largo del tiempo, los autores introducen una simplificación clave: asumen que la variación de la velocidad angular en el marco ECF es muy pequeña ($<3\%$) y, por lo tanto, puede tratarse como una constante ($\omega_{F}(t) \approx \omega_{F}$).

Bajo esta suposición, la distancia angular de la traza se linealiza respecto al tiempo mediante la \textbf{ecuación (10)}:
\begin{equation}
\psi(t) - \psi(t_{0}) \approx \omega_{F}(t - t_{0})
\end{equation}

Sustituyendo esta relación en la ecuación principal, la expresión que modela la \textbf{curva Doppler en función del tiempo} queda de la siguiente forma:

\begin{equation}
\frac{\Delta f}{f} \approx -\frac{1}{c}\frac{r_{E}r \sin(\omega_{F}(t-t_{0}))\cos(\gamma(t_{0}))\cdot\omega_{F}}{\sqrt{r_{E}^{2}+r^{2}-2r_{E}r \cos(\omega_{F}(t-t_{0}))\cos(\gamma(t_{0}))}}
\end{equation}

Donde el ángulo central constante en el momento de máxima elevación, $\gamma(t_{0})$, se define a partir de la \textbf{ecuación (4)}:
\begin{equation}
\cos(\theta_{max}+\gamma(t_{0}))=\frac{r_{E}}{r}\cos \theta_{max}
\end{equation}

Asimismo, la velocidad angular constante del sistema se aproxima mediante su valor en la latitud más alta (donde la velocidad tangencial relativa es mínima), expresada en la \textbf{ecuación (9)}:
\begin{equation}
\omega_{F} \approx \omega_{s} - \omega_{E}\cos i
\end{equation}

Donde $\omega_{s}$ es la velocidad angular del satélite en el marco inercial (ECI), $\omega_{E}$ es la velocidad angular de rotación de la Tierra e $i$ representa la inclinación de la órbita del satélite LEO.



\begin{table}[htbp]
\centering
\caption{Payload Downlink Budget for the CHESS Mission (X-band)}
\renewcommand{\arraystretch}{1.2}
\begin{tabular}{l l c c c l}
\toprule
\textbf{Item} & \textbf{Symbol} & \textbf{Adverse} & \textbf{Nominal} & \textbf{Favourable} & \textbf{Unit} \\
\midrule
Frequency & $f$ & 10.475 & 10.475 & 10.475 & GHz \\
Data Rate & $R$ & 10 & 25 & 25 & Mbps \\
S/C Output Power & $P$ & 0.47 & 1.48 & 5.29 & W \\
S/C Antenna Gain & $G_t$ & 17 & 19 & 21 & dBi \\
EIRP & - & 9.27 & 17 & 25.2 & dBW \\
S/C Max. Altitude & $H$ & 505 & 500 & 495 & km \\
GS Min. Elevation & $\varepsilon$ & 15 & 15 & 15 & deg \\
GS Antenna Gain & $G_r$ & 40 & 40.7 & 43.1 & dBi \\
GS Noise Temperature & $T_0$ & 148.5 & 135 & 121.5 & K \\
GS Figure of Merit & $G/T$ & 15.2 & 15.7 & 16.1 & dB/K \\
Various Losses\textsuperscript{a} & $L_{misc}$ & -6.23 & -5.47 & -4.33 & dB \\
\midrule
$E_b/N_0$ & - & 6.1 & 9.1 & 20.6 & dB \\
Required $E_b/N_0$ & - & 2.3 & 2.3 & 2.3 & dB \\
\textbf{Margin} & - & \textbf{3.8} & \textbf{6.8} & \textbf{18.3} & \textbf{dB} \\
\bottomrule
\multicolumn{6}{l}{\footnotesize \textsuperscript{a} Includes S/C \& GS pointing, line, atmospheric, and polarization losses.}
\end{tabular}
\end{table}


\begin{cpp}
#include <iostream>

int main() {
    // Este es un comentario en C++
    std::cout << "¡Hola desde C++!" << std::endl;
    return 0;
}
\end{cpp}